\documentclass[11pt]{article}
\ifdefined\pdftrailerid\pdftrailerid{}\fi
\usepackage[T1]{fontenc}
\usepackage{lmodern}
\usepackage[margin=.88in]{geometry}
\usepackage{amsmath,amssymb,amsthm,mathtools,microtype}
\usepackage{xcolor}
\usepackage[colorlinks,linkcolor=blue!45!black,citecolor=blue!45!black,urlcolor=blue!55!black]{hyperref}
\hypersetup{pdftitle={Adjoint torsion and a coupled logarithmic cutoff for paired-cap WRT invariants},pdfsubject={Quantum topology, adjoint torsion, and singular Chern--Simons sewing}}
\newtheorem{theorem}{Theorem}
\newtheorem{proposition}[theorem]{Proposition}
\newtheorem{lemma}[theorem]{Lemma}
\newtheorem{corollary}[theorem]{Corollary}
\newcommand{\SU}{\mathrm{SU}(2)}
\newcommand{\SO}{\mathrm{SO}(3)}
\newcommand{\RR}{\mathbb R}
\newcommand{\ZZ}{\mathbb Z}
\newcommand{\ii}{\mathrm i}
\newcommand{\Ad}{\operatorname{Ad}}
\newcommand{\sinc}{\operatorname{sinc}}
\newcommand{\vol}{\operatorname{Vol}}
\newcommand{\CS}{\operatorname{CS}}
\newcommand{\Haar}{\mathrm{Haar}}
\title{Adjoint torsion and a coupled logarithmic cutoff\\for paired-cap Witten--Reshetikhin--Turaev invariants}
\author{}
\date{}
\begin{document}
\maketitle
\begin{abstract}
For an unframed tree plumbing with \(+2,-2\) caps at each site, we compute the leading \(\SU\) Witten--Reshetikhin--Turaev measure with the fibre conjugacy angles on one side of a minimum cover held away from the central values. The limit is the Haar convolution of the actual cap holonomies. Computing the adjoint determinant lines gives the same radial measure, up to a constant, from square-root Reidemeister torsion. This comparison holds for every minimum cover: a seam with both fibres central contributes a rotation gluing parameter rather than another radial or determinant factor. For the four-site path with its middle sites as the cover, an independent radial cutoff of the true torsion gives a square of logarithmic scales, with twice the exact quantum weight. The actual reduced Chern--Simons quartic instead resolves the two normal directions only when the product of the two exterior fibre deficits exceeds inverse level on the logarithmic scale. Projected classical marks and literal finite-level Wilson marks then give the same triangle in the physical radial exponents. This determines the coupled cutoff without assuming a globally normalized Chern--Simons integration measure.
\end{abstract}

\section{Introduction and conventions}
In the semiclassical approximation to Chern--Simons theory, square-root adjoint torsion is the density at a smooth flat critical family; its role has been tested in lens spaces, torus bundles and finite-order mapping tori \cite{Jeffrey,AH,Charles}. At a singular endpoint this density can have a nonintegrable pole. Even if the density on the smooth branches is known, an independent cutoff on each endpoint variable need not reproduce the leading quantum coefficient. We give a family where the smooth density can be computed on all minimum-cover charts, and a closed example where the missing coupled cutoff is visible in both the actual Chern--Simons action and the exact modular sum.

Let \(G=(V,E)\) be a connected finite tree with \(N\ge2\) vertices. At each vertex take the oriented Seifert block over a sphere with \(\deg(v)\) outer boundary circles and two order-two exceptional fibres, obtained by opposite \(+2,-2\) fillings. Sew the outer tori by the marked fibre--section exchange along the edges of \(G\), with zero backbone framing. This is the usual Hopf-plumbing convention of \cite{GPPV,P48}. Write the closed manifold as \(M_G\). We use shifted level \(r=k+2\), the sine matrix \(S_{ab}=\sqrt{2/r}\sin(\pi ab/r)\), and \(\tau_r(S^3)=1\); the alternative convention is \(Z_r=S_{11}\tau_r\), so \(Z_r(S^2\times S^1)=1\). Denote the block fibre by \(h_v\). In the matched graph marking the two cap holonomies satisfy \(U_v^2=h_v^{-1}\), \(V_v^2=h_v\); their product is the inverse of the ordered product of section holonomies. Section holonomies are the adjacent fibres, with the appropriate orientations. Central signs in these products will not affect the Haar convolutions. The spectrum of the fundamental fibre Wilson operator is \(2\cos(\pi a/r)\), \(1\le a<r\). Its negative-central deficit is the actual Wilson coordinate \(x=(r-a)/r\), with classical counterpart \(x=1-\arg(h)/\pi\), where \(\arg(h)\in[0,\pi]\) is the \(\SU\) conjugacy angle. At finite level, a radial mark can equivalently be inserted by polynomial interpolation on these distinct eigenvalues.

Fix a minimum vertex cover \(A\), with independent complement \(B\). Put \(\nu=|A|\), \(\alpha=|B|=N-\nu\). A maximum matching pairs every \(a=j\in A\) with a distinct \(b_j\in B\): each of its \(\nu\) edges must use a cover vertex, hence no matched edge can use two cover vertices. Let \(B_*\) denote the unmatched vertices in \(B\). For the regular branch studied here, \(h_A=-1\), \(h_B\) is noncentral, and all \(A\) block representations are irreducible. At a given \(A\) site, reducibility can occur only when a signed sum of its noncentral boundary angles and the two cap angles \(\pi/2\) is a multiple of \(2\pi\) up to a fixed central shift. Thus away from finitely many proper angular hyperplanes all these block representations are irreducible. In particular the full stabilizer is the finite centre. Haar measures on compact groups and their conjugacy classes below are probabilities unless specified otherwise.

Our comparison separates three limits: the fixed-radius quantum limit and the determinant-line calculation (Sections~\ref{sec:quantum}--\ref{sec:torsion}), the independent radial regulation of that measure (Section~\ref{sec:regulator}), and the joint singular limit that retains the middle seam of the four-site path (Section~\ref{sec:p4}). The last result uses a particular exact analytic CS normal form whose general proofs and extensions are treated in the companion \cite{CSNF}. Earlier closed asymptotics \cite{P48} and cap phases and degeneration results \cite{P49} are inputs or comparators, not new claims here.

\section{The macroscopic WRT measure}\label{sec:quantum}

Define \(d\nu_0(x)=2\cos^2(\pi x/2)\,dx\) on \((0,1)\). For each \(b\in B\), choose \(h_b\) independently on its Haar conjugacy class of angle \(\pi(1-x_b)\), and let \(G_a\) be the ordered product of these \(h_b\)'s or their inverses at \(a\in A\) (central factors may be discarded). Set
\begin{equation}\label{eq:haar}
 K(g)=\frac{\pi}{4\sin(\arg g)},\qquad L_a(x)=\mathbb E K(G_a),\qquad
 dM_B(x)=\prod_{b\in B}d\nu_0(x_b)\prod_{a\in A}L_a(x).
\end{equation}
Singular values are interpreted as positive extended densities; the regularity and integrability used below follow from the proof.

\begin{theorem}[Macroscopic charge chart]\label{thm:macro}
At every \(a\in A\) insert a smooth deficit cutoff equal to one near \(0\) and supported in \([0,\epsilon)\), where \(\epsilon\max_v\deg(v)<1/4\). At every \(b\in B\) insert a test \(\psi_b\in C_c^\infty((0,1))\) in its fibre Wilson deficit. Then along all integer levels
\[
 \frac{\tau_r^{\rm mark}(M_G)}
 {(-1)^{r(N-1)}\,2^{3N/2-2\alpha}\pi^{-N}\,r^{N/2+\alpha}}
 \ \longrightarrow\ \int\prod_b\psi_b(x_b)\,dM_B(x).
\]
The measure in \eqref{eq:haar} is also a direct compact-group sewing measure on the specified flat branch.
\end{theorem}

\begin{proof}
In the classical branch, the sections of \(b\in B\) are central. Draw \(U_b\) by \(\SU\) Haar, then set \(h_b=U_b^{-2}\) and \(V_b\) by the product equation. Both cap squares hold; the Weyl density \((2/\pi)\sin^2\theta\,d\theta\) for \(\arg U_b=\theta\) pushes forward under \(x=|2\theta/\pi-1|\) to \(d\nu_0\), and its axis is uniform conditional on its class. At \(a\in A\) the caps are traceless. Their two independently Haar distributed unit imaginary quaternions have a product with uniform scalar part on \([-1,1]\). Its angle density \(\tfrac12\sin\phi\,d\phi\), divided by Haar's \((2/\pi)\sin^2\phi\,d\phi\), is exactly \(K\). Sewing the section product with a Haar delta therefore gives \(K(G_a)\).

For any bounded class factors, their joint angular average is the product of their averages. Indeed the bipartite incidence graph on \(A,B\), with the \(A\)--\(A\) edges deleted, is a forest. There is a factor with at most one \(B\) variable shared with the remaining factors (prune degree-one \(B\) vertices first). Integrate its private \(B\) axes. By conjugation invariance the resulting conditional expectation depends only on the fixed conjugacy class of the remaining shared \(B\) variable. Remove this factor and iterate. Nonnegative approximation extends this to \(K\).

For the quantum calculation put \(n_v=r-a_v\) and \(F_d(x)=\cos^2(\pi x/2)(\pi x/\sin\pi x)^d\). The tree surgery sum with the opposite caps is \cite{GPPV,P48}
\begin{equation}\label{eq:capped}
 \tau_r=(-1)^{r(N-1)}
 \bigl(2^{3N/2}\pi^{-N}+o(1)\bigr)r^{N/2}
 \sum_{\substack{1\le n_v<r\\ n_v\ {\rm odd}}}
  \prod_v F_{\deg(v)}(n_v/r)\prod_{vw\in E}\sinc(\pi n_v n_w/r).
\end{equation}
Here and below \(\sinc t=\sin t/t\) and the common scalar can be used unchanged for marked sums. To see the factors, the two cap amplitude is \(\left|(S\Theta^2S)_{r-n,1}\right|^2=(4/r)\mathbf1_{n\ {\rm odd}}\cos^2(\pi n/(2r))\), with \(\Theta_{aa}=e^{\pi\ii(a^2-1)/(2r)}\). For example the sine-product expansion reduces to two quadratic Gauss sums with linear charges differing by two; shifting the variable relates them by \(e^{-\pi\ii a/r}\). Squaring the Gauss sum leaves only differences \(0,r\) on the \(2r\)-grid, giving the parity projection and the stated sine square in \(a/r\). The vertex factor in \(Z_r\) is \(S_{1a_v}^{-\deg(v)}\) after the two terminal vertices combine into this amplitude. Factor \(\pi n_v n_w/r\) from each edge sine, leaving the sign \((-1)^r\) per edge. The paired caps cancel in the Schur-complement signature and the remaining tree adjacency has bipartite spectrum, so there is no residual surgery-signature phase in \eqref{eq:capped}. Divide \(Z_r\) by \(S_{11}\) to obtain the displayed scalar asymptotic.

The \(B\) sums are independent at fixed \(n_A\). Upon stripping \(r/4\) at each such vertex, the bracket is
\[
 U_{b,r}^{\psi}=\frac4r\sum_{\substack{l<r\\ l\ {\rm odd}}}
 \psi_b(l/r)\cos^2(\pi l/(2r))
 \prod_{a\sim b}\frac{\chi_{n_a}(\pi l/r)}{n_a},
 \qquad \chi_n(\theta)=\frac{\sin(n\theta)}{\sin\theta}.
\]
For any fixed \(n_A\) it tends to the corresponding \(d\nu_0\) integral; the \(A\)--\(A\) sincs and cutoffs tend to one. Uniform summability in \(n_A\) is essential here. Expand the numerator sines at each \(b\) into exponentials of signed frequencies \(q=\sum_{a\sim b}\pm n_a\). Repeated discrete summation by parts gives \(C_M(1+|q|)^{-M}\) per frequency after the \(4/r\) normalization. The amplitude is smooth with interior support after dividing by the powers of \(\sin(\pi l/r)\); its \(M\)-th grid difference is \(O_M(r^{-M})\) and \(|1-e^{2\pi\ii q/r}|\gtrsim |q|/r\) for all permitted \(|q|\le r/4\). In particular \(b_j\)'s bracket has rapid decay in \(n_j\) whenever \(n_j>2\sum_{i\ne j,\ i\sim b_j}n_i\). Starting at a largest \(A\) index, if this fails, follow a predecessor \(i\to j\) with \(n_i\ge n_j/(2D)\), \(D=\max\deg(v)\). These arrows are the cross edges after contracting the matching, hence their underlying graph is a forest. The walk must terminate at a dominant anchor comparable to \(\max n_A\) with graph-dependent constants. Thus the whole product of brackets decays faster than any power of \(\max n_A\), proving summable domination.

To identify the resulting continuum Fourier series, use the positive Abel kernels
\[
 K_t(g)=\sum_{\substack{n>0\\n\ {\rm odd}}}t^n\chi_n(\arg g)/n,\qquad 0<t<1.
\]
For \(\phi=\arg g\in(0,\pi)\),
\[
 K_t(g)=\frac{\arctan(2t\sin\phi/(1-t^2))}{2\sin\phi}
 \nearrow K(g).
\]
Schur averaging gives \(\mathbb E\chi_n(G_a)=n\prod_{b\sim a}\chi_n(\pi x_b)/n\) for odd \(n\); the initial factor \(n\) cancels the \(1/n\) in \(K_t\). Use the forest factorization and integrate the \(B\) variables, then send \(t\uparrow1\). The same rapid-decay estimate on the integrated coefficients justifies termwise limits, and positive monotone convergence identifies the density for nonnegative tests. Differences of nonnegative smooth tests cover the signed case. Combining the \(r/4\) factors with \eqref{eq:capped} proves the stated normalization.
\end{proof}

\section{The true adjoint torsion density}\label{sec:torsion}
We use positive cochain torsion \(T\), with the convention that an invertible \(0\)-to-\(1\) differential contributes the absolute determinant. Its acyclic analytic convention is \(\prod_q\det'(\Delta_q)^{(-1)^{q+1}q/2}\). For a closed irreducible flat local system in dimension three, take the tangent basis in \(H^1\) and its dual in \(H^2\) under Poincar\'e duality. Then \(T^{1/2}\) is a density in tangent coordinates: changing the \(H^1\) basis and its \(H^2\) dual contributes the square of the tangent determinant before taking the square root. We use the torsion gluing rule for a short exact sequence of based complexes, including the torsion of the long cohomology sequence \cite{Milnor}. The closed torsion has the equivalent Ray--Singer norm by Cheeger--M\"uller \cite{CM1,CM2}; no analytic gluing formula for pieces with boundary is needed.

\begin{theorem}[Determinant-line identification on every cover]\label{thm:torsion}
On the regular branch \(h_A=-1\), \(h_B\) noncentral, the pushforward of the actual positive adjoint torsion density \(T^{1/2}\) to the exterior radii \(x_B\), summing the regular bending and cap-root sheets, is \(C_{G,A}\,dM_B\). Here \(C_{G,A}>0\) depends on normalization and the cut, not on \(x_B\). The statement holds for a minimum cover with \(A\)--\(A\) seams as well as for an independent cover. The regular branch has dimension \(3N-2\nu-3\).
\end{theorem}

\begin{proof}
We compute the cochain determinants on the pieces, the long cohomology sequence on the cut tori, and finally the Poincar\'e pairing that turns the result into a tangent density.

\paragraph{\(B\) block, valence \(k\).}
Because its fibre \(h_b\) is noncentral, the adjoint system preserves an axial line and its perpendicular plane. All the block holonomies fix the axis; in particular the axial local system is trivial. Remove both exceptional solid tori. The remaining bundle is \(\Sigma_{0,k+2}\times S^1\); the plane circle differential \(\Ad(h_b)-I\) is invertible. Filter the product cochain complex by cellular base degree. Each quotient is the circle complex with that differential, with its torsion sign shifted by its cell degree. Their product gives plane torsion \(|\det(\Ad(h_b)-I)|^{\chi(\Sigma)}=|\det(\Ad(h_b)-I)|^{-k}\). The cap core is the section generator \(U_b\) or \(V_b\), up to inversion, in the actual primitive \((2,\pm1)\) filling slopes. Its solid torus contributes \(|\det(\Ad(U_b)-I)|\) or \(|\det(\Ad(V_b)-I)|\) to the plane torsion. Because the sections of \(B\) are central, \(V_b=\pm U_b^{-1}\). The plane on each boundary torus has identical nonzero determinants in successive degrees, hence torsion one and no cohomology. Here all determinants are restricted to the plane. The multiplicative rule therefore gives exactly
\begin{equation}\label{eq:B}
 T_{\rm plane}(B_b)^{1/2}
 =\frac{|\det(\Ad(U_b)-I)|}{|\det(\Ad(h_b)-I)|^{k/2}}
 =\frac{4\cos^2(\pi x_b/2)}{(2\sin(\pi x_b))^k}.
\end{equation}
The axial cohomology and determinants are constant in the marked slopes. Its \(H^0\) is \(\RR\). Its \(H^1\) restricts at the \(k\) tori to a common fibre value and independently a sum-zero section vector; the opposite exceptional slopes cancel in the base relation. Its \(H^2\) restricts to a sum-zero vector of oriented area classes. These descriptions also follow directly from ordinary cohomology: in real coefficients the two fillings express the exceptional section periods as opposite halves of the fibre period, and the outer boundary cycles generate the planar base homology subject to their sum.

\paragraph{\(A\) block.}
Its adjoint fibre holonomy is identity and its base representation is irreducible. After removing the exceptional solid tori, write \(W'=H^1(\Sigma;\Ad)\), where \(\Sigma\) includes all boundary circles. \(H^0(\Sigma;\Ad)=0\) and \(H^2(\Sigma;\Ad)=0\). The product has \(H^1=W'\), \(H^2=\eta\wedge W'\), where \(\eta\) evaluates as \(1\) on a regular fibre. In paired bases the two shifted copies of the surface complex have cancelling torsions. The boundary trace map from \(W'\) to the sum of the boundary invariant Lie algebras is onto by the relative sequence: \(H^2(\Sigma,\partial\Sigma;\Ad)=H^0(\Sigma;\Ad)^*=0\). These invariant algebras are axial lines at the caps and the \(A\)--\(B\) sections, and \(\RR^3\) at \(A\)--\(A\) sections.

For each removed cap, its core has adjoint rotation \(\pi\), so all nonzero plane torsion factors are constant. In the axial cohomology sequence the \(H^0\) of the solid torus maps isomorphically to that of its boundary. Its \(H^1\) maps with nonzero fibre coefficient of magnitude two in the marked slope bases, while the surface deformation maps to the independent section direction. The resulting degree-one differential onto the seam cohomology is onto. The degree-two restriction from the unfilled piece maps \(\eta\wedge W'\) onto the cap trace lines. These cohomology determinants are fixed slope determinants in splitting bases for the surjective boundary trace map. Thus the filled \(A\) piece has \(H^1=W_a\), \(H^2=\eta\wedge W_a\), \(H^0=H^3=0\), where \(W_a\subset W'\) has both cap traces fixed. The positive torsion of the piece is constant in paired bases. The real class \(\eta\) may be adjusted to extend across the caps by splitting the ordinary first cohomology with fibre value \(1\); the base has boundary. Its horizontal adjustment does not change the stated \(H^2\) description because absolute \(H^2(\Sigma;\Ad)=0\).

Write \(W_a^0\) for the kernel of the remaining seam restriction. These are the fixed-class tangent modes of the punctured sphere with its two cap classes and its noncentral boundary classes fixed; a central boundary class is a point and its deformation restricts by zero in the full \(\RR^3\). If \(k_a\) is the number of \(B\) neighbours, \(\dim W_a^0=2k_a-2\): there are \(k_a+2\) two-dimensional conjugacy classes, and the regular product constraint and effective conjugacy quotient each subtract three. Let \(\Omega_a\) be the Goldman symplectic form on that fixed-class space.

\paragraph{Seams and the closed long exact sequence.}
At an \(A\)--\(B\) seam only the axial cohomology survives, with dimensions \(1,2,1\). The plane is acyclic with torsion one. In oriented coordinates \((A\text{ fibre},A\text{ section})=(B\text{ section},B\text{ fibre})\), the \(A\) degree-one image is in the section direction, and the \(B\) image has sum-zero \(A\)-fibre coordinates incident to that \(B\), together with a common \(A\)-section value. At an \(A\)--\(A\) seam both \(\SU\) holonomies are central. Its adjoint torus complex is ordinary cohomology tensored with \(\RR^3\). The two \(A\) degree-one section maps go to complementary three-dimensional summands under the swap. They are onto their summands simultaneously with the other boundary maps by the surjectivity above. Matching therefore forces both section restrictions to vanish and yields no new degree-one cokernel at that seam. The \(H^0(T_e)=\RR^3\) survives in the gluing tangent space, and the difference of the two \(A\) degree-two restrictions is onto \(H^2(T_e)=\RR^3\); its kernel keeps one common \(\RR^3\) datum.

In more detail, let \(E_{AB},E_{AA}\) be the two edge sets. The global \(H^0,H^3\) vanish. The connecting image of the seam \(H^0\)'s in \(H^1(M_G;\Ad)\) is
\[
 L=\bigl(\RR^{E_{AB}}/\text{one diagonal per }b\in B\bigr)
       \ \oplus\ \bigoplus_{e\in E_{AA}}\RR^3 .
\]
The remaining \(H^1\) is the sum of the \(B\) radii and \(\oplus_a W_a^0\): matching forces all the \(B\) section vectors to be zero and the \(A\) section traces at \(A\)--\(B\) seams to be their \(B\) fibre periods. The connecting image of the seam \(H^1\)'s in \(H^2(M_G;\Ad)\) is one axial quotient per \(b\), given by the \(A\)-fibre vectors modulo the \(B\) sum-zero vectors. The remaining \(H^2\) consists of \(\oplus_a W_a^0\) in their \(\eta\)-copies, one \(B\) sum-zero area vector at every \(B\), and the common degree-two \(\RR^3\) at every \(A\)--\(A\) seam. To obtain these descriptions choose the trace splitting bases at \(A\) using the onto restriction map. The whole cohomology correction has constant matrices in these bases. Rotating between orthonormal adjoint frames to use the actual seam identifications changes their absolute determinants by one. In particular the number of tangent directions is
\[
 |B|+(|E_{AB}|-|B|)+3|E_{AA}|+\sum_{a\in A}(2k_a-2)
  =3N-2\nu-3.
\]

\paragraph{Dualizing the determinant line.}
The identity pairing a Mayer--Vietoris connecting class by restriction and cup product on the cut torus gives a triangular description of the closed Poincar\'e pairing. The axial degree-two connecting quotient pairs with the \(B\) fibre radii (their derivatives in the unit axis are \(\pm\pi\,dx_b\)); it is zero on the fixed-class and bending modes. The bending modes represented by \(L\) pair with the sum-zero axial degree-two data and with the common \(\RR^3\)'s, respectively, by constant quotient pairings. For the remaining pairing, use relative representatives of \(W_a^0\) with zero boundary trace and extend them by zero from that block. Pairing with \(\eta\) times the second representative integrates \(\eta\) to one along the regular fibre, leaving the fixed-class relative surface cup product \(\Omega_a\). The representatives can be chosen away from the cap neighbourhoods; no term from a cap or another block changes this middle pairing. Unspecified off-diagonal pairings do not change the triangular determinant. Replacing the \(H^2\) bases with Poincar\'e-dual ones therefore multiplies the positive torsion by \(\prod_a|\det\Omega_a|\) up to a constant. The square root supplies \(\prod_a|\operatorname{pf}\Omega_a|\), the surface Liouville measures.

The bending modes are actual gluing parameters. At each \(A\)--\(B\) seam their centralizer is \(U(1)\), with a common \(U(1)\) quotient at each \(B\); at an \(A\)--\(A\) seam it is \(\SU\) up to the finite centres. The tangent bases used above are unit-axis or orthonormal adjoint bases for these groups. Their volumes and finite quotients are independent of the radii in the regular chambers. A tree has no further loop condition on these identifications: fixing the piece representations at the prescribed classes leaves the indicated choices of edge intertwiners, modulo the vertex stabilizers. The two cap-root sheets per \(B\) also have fixed multiplicity.

\paragraph{Cancellation of the seam sine factors.}
For a fixed-class \(\SU\) conjugacy class of angle \(\theta\), the quasi-Hamiltonian Liouville measure is proportional to \(|\det(\Ad g-I)_{\perp}|^{1/2}\) times its normalized Haar measure, hence proportional to \(\sin\theta\) for noncentral classes, with a common constant \cite[Prop.~3.6]{AMW}. A central class is a point. The pushforward under fusion is convolution, and its value at the regular product level is proportional to the reduced symplectic volume \cite[\S5.1]{AMW}. These results apply at the regular radii specified above, since all configurations at the product level are irreducible there. The two cap sine factors are constants. Thus the volume at \(a\) is proportional to
\[
 \prod_{b\sim a,\ b\in B}\sin(\pi x_b)\ L_a(x).
\]
One can equally obtain this factor from the usual character formula for punctured-sphere volumes; at degree one the sum is understood in the heat or Abel sense. A central product shift and inversions are immaterial since \(K(-g)=K(g^{-1})=K(g)\). Each cross seam supplies exactly one sine factor. Together they cancel every sine denominator in \eqref{eq:B}. The remaining density is a constant times \(\prod_b2\cos^2(\pi x_b/2)\,dx_b\prod_aL_a(x)\), which is \eqref{eq:haar}. This proves the assertion, including the central--central case.
\end{proof}

\paragraph{Two-site check.}
The smallest path has two Klein-bottle interval bundles sewn by the fibre--section swap. On the irreducible branch with the left fibre central and the right one regular, the adjoint system splits globally into the right axis and its plane. In the Klein presentation \(XYX^{-1}=Y^{-1}\), the plane Fox differentials are \(d_0(v)=((X-I)v,(Y-I)v)\) and \(d_1(a,b)=(I-Y^{-1})a+(X+Y^{-1})b\). The left plane and the seam torus have torsion one by cancellation of the determinants with noncentral \(Y\). On the right \(Y=I\), so the plane square-root torsion is \(\sqrt{|\det(X-I)|/|\det(X+I)|}=\cot(\pi x/2)\). The axial matrices and the \(dx\) duality normalization are constant. This directly checks both endpoint behaviours in \eqref{eq:haar}, which here gives \(\frac{\pi}{4}\cot(\pi x/2)\,dx\).

\section{An independent endpoint regulator}\label{sec:regulator}

Put \(\beta(x)=\min(x,1-x)\). In \eqref{eq:haar} one can estimate each angular factor uniformly even when the incident radii are very unequal. Condition on all axes except that of one of the incident holonomies. Strip central signs so that its angle \(y=\pi\beta(x_b)\) and the angle \(z\) of the remaining product lie in \([0,\pi/2]\). The product has cosine \(\cos y\cos z-\sin y\sin z\,s\), where \(s\) is uniform on \([-1,1]\). Integration in this variable gives
\begin{equation}\label{eq:shell}
 \mathbb E\!\left[\frac1{\sin(\arg G_a)}\,\middle|\,\text{other axes}\right]
 =\frac{y+z-|y-z|}{2\sin y\sin z}
 =\frac{\min(y,z)}{\sin y\sin z}.
\end{equation}
In particular \(L_a\le C/\max_{b\sim a,\ b\in B}\beta(x_b)\). If \(x_{b_j}\to0\) and the other incident radii are \(o(x_{b_j})\), then \(x_{b_j}L_j\to1/4\). At any fixed noncentral anchor this conditional formula also gives continuity as other angles move or become central. The same estimates apply in all endpoint branches, with \(\beta\) for the distances.

Define \(i\to j\) on \(A\) when \(i\ne j\) is adjacent to \(b_j\). Let \(F\) consist of the ancestors (including the sites themselves) of the neighbours of unmatched vertices in \(B_*\), and \(D=A\setminus F\). The arrows have no cycles. Set \(P_{\rm ord}=\{s\in[0,1]^D:s_i\le s_j\text{ on }i\to j\}\), and \(d=|D|\). For forced \(A\) indices, write \(Q_b(n)=\Pr(\sum_{a\sim b}q_a=0)\), where \(q_a\) is uniform among the integers between \(-(n_a-1)/2\) and \((n_a-1)/2\), independently, with odd positive \(n_a\). Only forced neighbours are used at \(b\in B_*\) or \(b_j\) with \(j\in F\). Define
\[
 K_G=\sum_{(n_a)_{a\in F}}\prod_{b\in B_*}Q_b(n)\prod_{j\in F}Q_{b_j}(n).
\]
This is the bounded-defect factor of the exact paired-tree leading asymptotic in \cite{P48}. Empty products and zero-dimensional volumes are one.

\begin{proposition}\label{prop:reg}
If all \(B\) radii are cut off independently by \(\beta(x_b)\ge\delta\), then
\[
 M_B\{\beta(x_b)\ge\delta\ \forall b\}
 \ \sim\ 2^{-d}\vol(P_{\rm ord})K_G\bigl(\log(1/\delta)\bigr)^d
 \qquad(\delta\downarrow0).
\]
For an independent cover, \(P_{\rm ord}\) is the entire free order-polytope face of the exact WRT theorem. In a cover with \(A\)--\(A\) seams it need not be: the exact face imposes \(s_i+s_j\le1\) on those free seams in addition.
\end{proposition}
\begin{proof}
Write \(d_j=\beta(x_{b_j})\) and move \(d_j\) into the \(j\)-th factor \(L_j\) when rewriting its radial integration as \(dd_j/d_j\). Put \(\rho_j=d_j/\max_{b\sim j}\beta(x_b)\le1\). Each forced \(j\) has a directed path to a vertex adjacent to an unmatched \(b_*\). Along it the product of \(\rho\)'s telescopes to at most \(d_j/\beta(x_{b_*})\). Using these bounds and the trivial bound \(1\) gives
\[
 \prod_{j\in F}\rho_j\le
 \prod_{j\in F}d_j^{1/(2|F|)}
 \prod_{j\in F}\beta(x_{b_*(j)})^{-1/(2|F|)}
\]
when \(F\ne\varnothing\). This is integrable with the forced \(dd_j/d_j\) and unmatched ordinary measures. Keep those coordinates fixed and rescale the free logarithms by \(-\log d_j/\log(1/\delta)\). No free \(j\) is adjacent to an unmatched or a forced partner \(b_i\); otherwise it is forced. Outside the tie planes, \(d_jL_j\) tends to \(1/4\) when its own partner dominates all competitors, and to zero otherwise, by \eqref{eq:shell}. Simultaneously these conditions are precisely \(s_i<s_j\) on the free arrows \(i\to j\). The lower endpoint of the radial cap density tends to \(2\); the upper endpoint tends to zero and does not give a logarithmic coefficient. The uniform forced bound and bounded normalized free integrands give dominated convergence. Setting the free radii to their central limits in the forced factor is legitimate by \eqref{eq:shell} with their fixed noncentral anchors.

The remaining forced integral is \(K_G\). Indeed apply \(K_t\) as in the proof of Theorem~\ref{thm:macro}, and integrate each forced or unmatched \(b\) over \(0<x<1\). The normalized product of odd-dimensional characters is a trigonometric polynomial with only frequencies \(2\pi\) times integers. Integrating against \(2\cos^2(\pi x/2)=1+\cos(\pi x)\) retains just its constant coefficient, \(Q_b(n)\) (the product is a cosine polynomial). All these integrated coefficients are nonnegative, so \(t\uparrow1\) preserves their sum. Finite integrability above proves convergence. Each surviving free factor is \(1/2\). This gives the stated result. The exact WRT comparison follows from the independent-set face rule: saturating the matching gives \(t_{b_j}=1-t_j\), and an edge within \(A\) is still a sinc with the product \(r^{t_i+t_j-1}\). In the smooth chart regulator the coordinate \(s_j\) is the corresponding complement of the \(B\) charge exponent.
\end{proof}

\paragraph{Power check independent of the modular sum.}
Near the all-negative fibre locus with nonparallel adjacent axes, let \(S\) be the vertices with nonzero small signed fibre defects \(p_v\) in \(h_v=-\exp(-2p_vu_v)\). Commutation forces \(S\) to be independent. At each exterior neighbour \(a\in N(S)\), the other cap square is the one scalar constraint
\[
 u_a\cdot\log\!\left(\prod_{v\sim a}\exp(-2\epsilon_{av}p_vu_v)\right)=0 .
\]
The product order is the actual section order. Tree edge cosines are independent local coordinates off parallelity. Divide each row by the largest incident \(|p_v|\) and use its edge cosine as pivot. The linear-in-defect part gives a diagonal invertible selected derivative matrix; higher terms have \(O(\max|p|)\) derivative after division, since they each contain at least two defects from that row. Thus these are exactly \(|N(S)|\) independent equations in the small nonparallel chart. The stratum has dimension \(2N-3+|S|-|N(S)|\) where nonempty. It is at most \(3N-2\nu-3\): a largest independent set in the remaining forest has at least half of its \(N-|S|-|N(S)|\) vertices and can be added to \(S\). Equality is attained at \(S=B\) near a point with all cross cosines zero, by the same independent pivots and the implicit function theorem. Half the true maximum dimension agrees with the exact leading \(Z_r\) exponent \(N/2+\alpha-3/2\) in \cite{P48}. This is a stratum comparison within the stated chart, not an exclusion of other components.

\section{The four-site path and its coupled logarithmic cutoff}\label{sec:p4}
Let the vertices be \(B_1-A_1-A_2-B_2\), with the middle sites as cover. The positive macroscopic density in Theorem~\ref{thm:macro} is
\begin{equation}\label{eq:cot}
 dM_B=\left(\frac{\pi}{4}\right)^2
        \cot(\pi x_1/2)\cot(\pi x_2/2)\,dx_1\,dx_2 .
\end{equation}
Indeed each middle site has just one noncentral \(B\) neighbour. Theorem~\ref{thm:torsion} identifies the same shape from the true torsion, including the middle seam's \(\SO\) bending. Its independent cutoff \(x_i\ge 1/r\) near the negative endpoint gives \(\tfrac14(\log r)^2\). We now compare this to the actual off-shell action and the exact quantum joint limit in the same physical branch radii.

\subsection{The physical branch in the exact quartic chart}
Choose an all-negative fibre reference with each exterior pair of axes orthogonal and the middle cosine \(b\) with \(0<|b|<1\). The chamber \(h_{A_i}=-1\), both outer fibres regular near \(-1\) closes onto this locus. A signed exterior holonomy defect is \(p_{B_i}\) with \(h_{B_i}=-\exp(-2p_{B_i}u_{B_i})\). Its physical Wilson radius is \(x_i=(2/\pi)|p_{B_i}|\) near this endpoint.

\begin{lemma}[Exact local CS input and preservation of branch scales]\label{lem:cs}
In this reference chart the actual reduced \(\CS\) function has real analytic form
\begin{equation}\label{eq:quartic}
 \CS_0+\kappa s_1s_2 y_1y_2,\qquad \kappa\ne0,
\end{equation}
with three passive coordinates \(\phi\). The specified exterior-fibre branch is \(P=\{s_1=s_2=0\}\). On \(P\), the physical radius \(x_i\) is \(|y_i|\) times a positive analytic unit. Thus the logarithmic exponents of the two physical radii and of \(|y_i|\) differ by \(O(1/\log\lambda)\), uniformly even when their sizes are very unequal.
\end{lemma}
\begin{proof}
The analytic normal-form construction is proved in the companion \cite{CSNF}. We recall the specific steps giving the needed preservation of the physical branch. In local holonomy variables take \(U_v=u_v\exp(p_vu_v)\) and \(h_v=U_v^{-2}\). Nonparallelity makes the exact commuting ideal \(J_0=(p_vp_w:vw\in E)\). Modulo it, the other cap square on a block is the component along its axis of the logarithm of the ordered section product; conjugation by \(U_v\) reduces to the half-turn about \(u_v\). Write that scalar as \(\sum_w p_w(B_{vw}+c_{vw})\), where \(B_{vw}\) are the signed edge cosines and \(c_{vw}\) vanish to first order in the neighbouring defects. The symmetric correction \(T_{vw}=B_{vw}+c_{vw}+c_{wv}\) changes this equation only by elements of \(J_0\). Since these edge angles are independent local tree coordinates, the flat analytic equations are exactly the Jacobian equations of \(\sum_{vw}T_{vw}p_vp_w\). This works as an identity of analytic ideals, not only on real points.

In the \(P_4\) case put \(x_A=(p_{A_1},p_{A_2})\), \(y=(p_{B_1},p_{B_2})\), and write \(d_i\) for the corrected exterior dot. The graph function is \(\tfrac12 x_A^{\mathsf T}B x_A+x_A^{\mathsf T}(d\odot y)\), with \(B\) the invertible two-by-two adjacency of the middle edge and \(B_{12}\) nonzero near \(b\). Eliminating its split quadratic part leaves \(-d_1y_1d_2y_2/B_{12}\). An analytic unit change of the \(d\)'s freezes the denominator and leaves \(y\) intact. This is an elimination of the graph equations to compare critical ideals; it does not add two variables to the actual harmonic \(\CS\) chart.

For clarity, the bridge from those equations to the function is as follows. On a Coulomb slice at this irreducible flat connection the coexact transverse \(*d_A\) is invertible, so analytic stationary reduction leaves a finite function on harmonic one-forms. Its critical equation is flatness, also over local Artin algebras: once the coexact and harmonic projections of \(*F\) vanish, Bianchi and invertibility on zero-forms eliminate the exact part. Holonomy gives an analytic coordinate map from the reduced harmonic variables to the retained graph variables, invertible on tangent spaces by the infinitesimal holonomy isomorphism. Over each Artin algebra, flat holonomy and the unique local Coulomb representative are inverse descriptions \cite[\S6]{GM}. Equality of the ensuing completed Jacobian ideals descends to the convergent local ring. In particular the actual reduced function \(f\), with its constant removed, has \(J_f=J_q\), \(q=s_1s_2y_1y_2\). The cubic gradient span forces its leading jet to be a nonzero scalar \(\kappa q\): every degree-four monomial whose derivatives belong to that span contains all four active coordinates, each to degree one, and no spectator. Terms of lower degree vanish. Finally \(\nabla(f-\kappa q)=A(z)\nabla(\kappa q)\) with \(A(0)=0\), subtracting any constant homogeneous syzygy. Radial Euler integration gives \(f-\kappa q\in\mathfrak m^2J_q\). Interpolating between \(\kappa q\) and \(f\), the gradients remain equivalent by \(I+tA(z)\); the corresponding analytic Moser vector field is \(O(|z|^2)\). Its flow gives an actual right equivalence tangent to the identity. This is the specialized quartic assertion used here.

Before this final equivalence \(y_i=p_{B_i}\), and \(P\) represents the two noncentral exterior fibres with the middle fibres fixed and their exterior orthogonality conditions imposed. The quartic critical set is the union of six coordinate planes obtained by setting two of the four active coordinates to zero, times the passive space. The tangent-identity Moser equivalence preserves that set and fixes each component as a germ, since their tangent planes are distinct. It also preserves their intersections. Hence on \(P\) each of \(y_i=0\) is preserved, and the physical \(p_{B_i}\) is divisible on \(P\) by \(y_i\), with nonzero unit derivative. An earlier constant scale change can be kept in \(s\). This gives the asserted physical radii and the uniformly bounded log-scale error.
\end{proof}

\subsection{Conditional normal integration and a projected physical mark}
Let \(a(s,y,\phi)\) be any smooth compactly supported amplitude in a sufficiently small local quartic chart with \(|y_i|<1\) and physical \(x_i<1\). For a bounded continuous test \(F\) of two nonnegative logarithmic exponents, take its mark in the physical \(x_i\) on \(P\), extended to the chart by projection \((s,y,\phi)\mapsto(0,y,\phi)\). Use the phase \(e^{\ii\lambda(\CS-\CS_0)}\), and denote the resulting integral by \(I_\lambda[F]\).

\begin{proposition}\label{prop:classicalmark}
Set \(\Delta=\{(u_1,u_2):u_i\ge0,\ u_1+u_2\le1\}\). Then
\begin{equation}\label{eq:class}
 \frac{\lambda I_\lambda[F]}{(\log\lambda)^2}
 \longrightarrow
 \frac{8\pi}{|\kappa|}
  \left[\int a(0,0,0,0,\phi)\,d\phi\right]
  \int_\Delta F(u_1,u_2)\,du_1\,du_2 .
\end{equation}
For a single \(y\)-sign quadrant the prefactor is \(2\pi/|\kappa|\). The physical crossover is \(\lambda x_1x_2\) of order one, rather than one independent cutoff per radius.
\end{proposition}
\begin{proof}
At fixed exterior radii integrate the two \(s\) variables first. For \(T\in\RR\),
\[
 J(T;y,\phi)=\int a(s,y,\phi)e^{\ii T s_1s_2}\,ds,\qquad
 J(T;y,\phi)=\frac{2\pi}{|T|}a(0,0,y,\phi)+O(|T|^{-2})
 \quad (|T|\ge1).
\]
The remainder is uniform: take the Fourier transform in \(s_2\), set \(w=Ts_1\), and expand the smooth amplitude at \(s_1=0\). The first derivative has a uniformly integrable weighted Schwartz bound in \(w\). Direct integration for bounded \(T\) completes \(|J|\le C/(1+|T|)\). With \(T=\lambda\kappa y_1y_2\), this reproduces the \(1/(|y_1y_2|)\) resolved pole for large \(|T|\), as in \eqref{eq:cot}; it stays bounded in the unresolved region. In each sign quadrant put \(|y_i|=\lambda^{-u_i}\). After division by \((\log\lambda)^2/\lambda\) the absolute bound is
\[
 C\min\{1,\exp[-(\log\lambda)(u_1+u_2-1)]\}.
\]
For \(\log\lambda\ge1\) this has an integrable majorant on \(u_i\ge0\), and the passive support is fixed compact. On the interior of \(\Delta\) the leading resolved term applies; on its exterior the normalized integrand tends to zero. Lemma~\ref{lem:cs} gives the same limiting logarithmic exponents for the projected physical marks. Dominated convergence yields \eqref{eq:class}.
\end{proof}

The projection in this statement matters. Right equivalence preserving a critical scheme does not determine an original observable off the branch. For example \(q\) and \(q+\epsilon(s_1s_2y_2)^2\) have the same Jacobian ideal and are related by \(y_1\mapsto y_1+\epsilon s_1s_2y_2\). This is the identity on \(P\), but its changed \(y_1\) coordinate is not divisible by \(y_1\) off \(P\). We make no assertion comparing unprojected off-shell holonomies in such charts.

\subsection{Literal finite Wilson marks and the exact triangle}
Use the actual quantum exterior deficits \(x_i=l_i/r\), with \(l_i\) positive odd. Keep smooth fixed cutoffs on the middle deficits \(n_i/r\) near \(0\). Insert \(F= f_1(u_1)f_2(u_2)\), \(f_i\in C_c^\infty((0,1))\), with \(u_i=1-\log l_i/\log r\). In \eqref{eq:capped}, each pendant bracket after pulling out \(r/(4n_i)\) is
\[
 \frac4\pi\sum_{l\ {\rm odd}<r}\frac{F_1(l/r)}l
     \sin(\pi n_i l/r)\,f_i(1-\log l/\log r).
\]
It is uniformly bounded and tends to \(f_i(t_i)\) uniformly on compact interior intervals of \(t_i=\log n_i/\log r\). To prove this, set \(v=n_i l/r\). On \(\delta<v<R\) the mesh is \(2n_i/r\), \(F_1(l/r)\to1\), and the mark tends to \(f_i(t_i)\). The small-\(v\) part is \(O(\delta)\) by the sine bound. The tail \(v>R\) is \(O(1/R)\) by geometric summation by parts: on \(l\ge L\) the weight involving \(F_1\), \(f_i\), and \(1/l\) has size and total variation \(O(1/L)\). Compact support in the mark puts \(l\le r^{1-\epsilon'}\) for some \(\epsilon'>0\), so there is no reflected upper endpoint. For a uniform bound outside the compact \(t_i\) range split at \(l\) of order \(r/n_i\) and apply the same two estimates. The limiting odd-grid integral is \((2/\pi)\int_0^\infty\sin(\pi v)\,dv/v=1\).

The middle edge retains \(\sinc(\pi n_1n_2/r)\). Its limit is one for \(t_1+t_2<1\) and zero for \(t_1+t_2>1\), with a uniform bound of one. The odd harmonic measures, divided by \((\log r)/2\), converge to Lebesgue on the logarithmic unit interval. The bounded pendant brackets control arbitrarily small strips about the inequalities and endpoints. Writing \(\mathcal N_r=(-1)^r(4/\pi^4)r^4\) for the four-site macroscopic normalization, we get the literal Wilson-mark law
\begin{equation}\label{eq:qmark}
 \frac{\tau_r^{\rm mark}}{\mathcal N_r(\log r)^2}
  \ \longrightarrow\ \frac14\int_\Delta F(u_1,u_2)\,du_1du_2 .
\end{equation}
\begin{samepage}A rapidly convergent tensor expansion of a smooth compact mark gives the same statement for non-product marks; the bracket bounds use only finitely many seminorms of the individual tests. The known exact unmarked result \cite{P48}, with both endpoints treated, gives the same full leading mass \(1/8\) in these units. In contrast an independent cutoff of the fixed macroscopic measure \eqref{eq:cot} gives \(\tfrac14\,du_1du_2\) on the entire unit square, of total mass \(1/4\). Thus taking a fixed macroscopic limit and independently truncating its torsion singularities does not commute with the joint quantum limit. The actual normal Hessian yields the coupled threshold \(\lambda x_1x_2\) in the same branch scales, not just a triangle of equal area in unrelated transform coordinates.

\paragraph{What the comparison normalizes.}
Fix a regular bending chart and one signed radial sheet. If the resolved normal contribution after removing \(\lambda^{-1}\) is calibrated to the actual local torsion density \(\rho_P\,dy\,d\phi\), then its normal amplitude trace is forced to be \(a|_{s=0}=|\kappa y_1y_2|\rho_P/(2\pi)\). This is smooth through the corner: \(x_i/|y_i|\) are units, \(x_i\cot(\pi x_i/2)\) is even analytic there, and the bending and coordinate Jacobians are smooth. Any smooth extension of this trace gives the same leading projected integral \eqref{eq:class}. This is a relative local calibration, not a derivation of the full quantum measure. The residual \(\CS\) density, the overall conversion from the local \(\lambda^{-1}(\log\lambda)^2\) to \(Z_r\) of order \(r^{5/2}(\log r)^2\), and the global accounting of other angular charts and overlapping strata are not computed by this argument. The result assumes a fixed nonzero, nonparallel middle cosine for the CS chart. The middle-normal Hessian used here has zero signature; absolute APS and framing phases are not inferred from this local cutoff.

\par\end{samepage}


\begin{thebibliography}{12}
\setlength{\itemsep}{2pt}
\bibitem{Jeffrey} L.~C.~Jeffrey, \emph{Chern--Simons--Witten invariants of lens spaces and torus bundles, and the semiclassical approximation}, Comm. Math. Phys. \textbf{147} (1992), 563--604.
\bibitem{AH} J.~E.~Andersen and B.~Himpel, \emph{The Witten--Reshetikhin--Turaev invariants of finite order mapping tori II}, Quantum Topol. \textbf{3} (2012), 377--421, \href{https://arxiv.org/abs/1107.1813}{arXiv:1107.1813}.
\bibitem{Charles} L.~Charles, \emph{On the Witten asymptotic conjecture for Seifert manifolds}, \href{https://arxiv.org/abs/1605.04124}{arXiv:1605.04124}.
\bibitem{GPPV} S.~Gukov, D.~Pei, P.~Putrov and C.~Vafa, \emph{BPS spectra and 3-manifold invariants}, J. Knot Theory Ramifications \textbf{29} (2020), 2040003, \href{https://arxiv.org/abs/1701.06567}{arXiv:1701.06567}, Appendix A.
\bibitem{AMW} A.~Alekseev, E.~Meinrenken and C.~Woodward, \emph{Duistermaat--Heckman measures and moduli spaces of flat bundles over surfaces}, Geom. Funct. Anal. \textbf{12} (2002), 1--31; \href{https://arxiv.org/abs/math/9903087v2}{arXiv:math/9903087v2}, listed as \emph{Duistermaat--Heckman distributions for group valued moment maps}.
\bibitem{CM1} J.~Cheeger, \emph{Analytic torsion and the heat equation}, Ann. of Math. \textbf{109} (1979), 259--321, \href{https://annals.math.princeton.edu/1979/109-2/p02}{publisher page}.
\bibitem{CM2} W.~M\"uller, \emph{Analytic torsion and R-torsion of Riemannian manifolds}, Adv. Math. \textbf{28} (1978), 233--305, \href{https://doi.org/10.1016/0001-8708(78)90116-0}{doi:10.1016/0001-8708(78)90116-0}.
\bibitem{Milnor} J.~Milnor, \emph{Whitehead torsion}, Bull. Amer. Math. Soc. \textbf{72} (1966), 358--426.
\bibitem{GM} W.~M.~Goldman and J.~J.~Millson, \emph{The deformation theory of representations of fundamental groups of compact K\"ahler manifolds}, Publ. Math. IH\'ES \textbf{67} (1988), 43--96, \href{https://numdam.org/item/PMIHES_1988__67__43_0/}{\S6 (flat holonomy)}.
\bibitem{P48} \emph{Leading logarithms in closed \(\SU\) Witten--Reshetikhin--Turaev invariants of capped trees}, research manuscript 48, \href{https://github.com/g3173028019-beep/ai-automated-quantum-gravity-research/tree/d23f0b109029d37a09a1b8758dd73f1e27a7ef72/papers/48-wrt-leading-logarithms}{release revision}.
\bibitem{P49} \emph{Chern--Simons phases and degenerate sewing in \(\SU\) WRT invariants}, research manuscript 49, \href{https://github.com/g3173028019-beep/ai-automated-quantum-gravity-research/tree/d23f0b109029d37a09a1b8758dd73f1e27a7ef72/papers/49-wrt-semiclassical-sewing}{release revision}.
\bibitem{CSNF} \emph{Exact Chern--Simons normal forms for paired-cap tree plumbings}, \href{https://github.com/g3173028019-beep/ai-automated-quantum-gravity-research/tree/main/papers/49b-cs-normal-forms}{companion manuscript}.
\end{thebibliography}
\end{document}
